%!TEX root = main.tex

% load necessary packages
\usepackage[utf8]{inputenc}
\usepackage[T1]{fontenc}
\usepackage{amsmath,,amssymb,amsthm,amsfonts,amsxtra,mathtools}
\usepackage{amsfonts}
\usepackage[margin =3.25cm]{geometry}
\usepackage{amssymb}
\usepackage{graphicx}
\usepackage[mathscr]{euscript}
\usepackage[english]{babel}
% \usepackage{cite}
\usepackage{lmodern}
\usepackage{appendix}
%Auflistungen ändern
\usepackage{enumerate}
\usepackage{titlesec}
\usepackage{titletoc}
\usepackage{fancyhdr} %header package
\usepackage{setspace}
\usepackage{todonotes}
\usepackage{hyperref}
\usepackage{adjustbox}
\usepackage{graphicx}

% section formating
\titleformat{\chapter}[display]
  {\normalfont\huge\bfseries}{}{0pt}{\Huge}
\renewcommand\thesection{\arabic{section}}
\renewcommand{\headrulewidth}{0pt}
\titlecontents{chapter}
  [0em]{\addvspace{10pt}\bfseries}
  {}{}{\titlerule*[1000pc]{.}\contentspage}

%set page style and define header
\pagestyle{fancy}
\fancyhf{}
\fancyhead[L]{\nouppercase{\rightmark}}
\fancyfoot[C]{\thepage}

%bibliography
\usepackage[
    backend=biber,
    natbib=true,
    style=alphabetic,
    mincitenames = 4,
    maxcitenames = 6,
    minbibnames = 4,
    maxbibnames = 6,
    sorting=nty,
    url = false,
    doi = false,
    isbn = false,
    eprint =false
]{biblatex}
\addbibresource{bibliography.bib}

%delete p &pp in before optional argument
% \DeclareFieldFormat{postnote}{#1}
% \DeclareFieldFormat{multipostnote}{#1}
% \DeclareFieldFormat[article, incollection]{title}{#1}

% hspace für bl;de Fehler bei Defs, etc
\newcommand{\h}{\hspace{-0.5mm}}

%nices lifting symbol
\usepackage{tikz}
  \newlength\squareheight
  \setlength\squareheight{3.5pt}
  \newcommand\squareslash{\tikz{\draw (0,0) rectangle (\squareheight,\squareheight);\draw(0,0) -- (\squareheight,\squareheight)}}
% \newcommand{\pr}{to [out=0,in=180]}
% \usetikzlibrary{math}
% \usetikzlibrary{decorations.markings}
% \usetikzlibrary{decorations.pathreplacing}
% \usetikzlibrary{arrows,shapes,positioning}
% %%multiple arrow options
% \tikzstyle directed=[postaction={decorate,decoration={markings,
%     mark=at position #1 with {\arrow{>}}}}]
% \tikzstyle rdirected=[postaction={decorate,decoration={markings,
%     mark=at position #1 with {\arrow{<}}}}]

% \tikzset{
%     anchorbase/.style={baseline={([yshift=-0.5ex]current bounding box.center)}},
%     tinynodes/.style={font=\tiny,text height=0.75ex,text depth=0.15ex},
%     smallnodes/.style={font=\scriptsize,text height=0.75ex,text depth=0.15ex},
%     >={Latex[length=1mm, width=1.5mm]},
%     rd/.style={red, line width=.5mm},
%     bl/.style={blue, line width=.5mm},
%     gr/.style={green, line width=.5mm},
%     pu/.style={purple, line width=.5mm},
%     or/.style={orange, line width=.5mm},
%     web/.style={very thick},
%     dweb/.style={double},
%     wh/.style={white,line width=.15cm}
%   }
% \tikzcdset{arrow style=tikz, diagrams={>=stealth}}
% \tikzcdset{scale cd/.style={every label/.append style={scale=#1},
%     cells={nodes={scale=#1}}}}


%Grafiken, Farben,...
\usepackage{xcolor}
\usepackage{tikz}
\usetikzlibrary{braids}
\usetikzlibrary{matrix}
\usepackage{tikz-cd}
\usepackage{subfigure}
\usepackage[arrow, matrix, curve]{xy}

%Style plain für Theoreme, Sätze, Lemma, Propositionen,...
\theoremstyle{plain}
\newtheorem{thm}{Theorem}[section]
\newtheorem{prop}[thm]{Proposition}
\newtheorem{lem}[thm]{Lemma}
\newtheorem{cor}[thm]{Corollary}
\newtheorem{cla}[thm]{Claim}

%Style definition für Definitionen und Beispiele...
\theoremstyle{definition}
\newtheorem{defin}[thm]{Definition}
\newtheorem{eg}[thm]{Example}
\newtheorem{egs}[thm]{Examples}
\newtheorem{rem}[thm]{Remark}
\newtheorem{nota}[thm]{Notation}

\newcommand{\R}{\mathbb{R}}
\newcommand{\N}{\mathbb{N}}
\newcommand{\Z}{\mathbb{Z}}
\newcommand{\Q}{\mathbb{Q}}
%Mathoperators
\DeclareMathOperator{\Set}{\textbf{Set}}
\DeclareMathOperator{\Plop}{\textbf{PlanarOp}}
\DeclareMathOperator{\Ch}{\textbf{Ch}}
\DeclareMathOperator{\Mod}{\textbf{-mod}}
\DeclareMathOperator{\Grp}{\textbf{Grp}}
\DeclareMathOperator{\Top}{\textbf{Top}}
\DeclareMathOperator{\Cat}{\textbf{Cat}}
\DeclareMathOperator{\Vect}{\textbf{Vect}}
\DeclareMathOperator{\id}{Id}
\DeclareMathOperator{\op}{op}
\DeclareMathOperator*{\colim}{colim}
\DeclareMathOperator{\Mul}{Mul}
\DeclareMathOperator{\Map}{Map}
\DeclareMathOperator{\Hom}{Hom}
\DeclareMathOperator{\HHom}{\underline{Hom}}
\DeclareMathOperator{\Emb}{Emb}
\DeclareMathOperator{\End}{End}
\DeclareMathOperator{\Ev}{Ev}
\DeclareMathOperator{\Cyl}{Cyl}
\DeclareMathOperator{\Path}{Path}
\DeclareMathOperator{\Fun}{Fun}
\DeclareMathOperator{\Ass}{Ass^{\otimes}}
\DeclareMathOperator{\dd}{D}
\DeclareMathOperator{\E}{E}
\DeclareMathOperator{\Ex}{\E_2^{\otimes}}
\DeclareMathOperator{\Comm}{Comm^{\otimes}}
\DeclareMathOperator{\Ho}{Ho}
\DeclareMathOperator{\coker}{coker}
\DeclareMathOperator{\T}{T}
\DeclareMathOperator{\I}{I}
\DeclareMathOperator{\inter}{int}




%nice Buchstaben fuer categorien
\usepackage[mathscr]{euscript}
% \renewcommand{\S}{\mathscr{S}}
\renewcommand{\P}{\mathscr{P}}
\renewcommand{\S}{\mathscr{S}}
\newcommand{\C}{\mathscr{C}}
\renewcommand{\O}{\mathscr{O}}
\newcommand{\B}{\mathscr{B}}
\newcommand{\D}{\mathscr{D}}
\newcommand{\nerv}{\mathscr{N}}
\newcommand{\W}{\mathscr{W}}
\newcommand{\Fib}{\mathscr{F} \text{ib}}
\newcommand{\Cof}{\mathscr{C} \text{of}}
\newcommand{\M}{\mathscr{M}}
\newcommand{\X}{\mathcal{X}}
\newcommand{\ind}{\text{ind-}\C}
% \newcommand{\I}{\mathcal{I}}
\newcommand{\chains}{\Ch_k^\I}
\newcommand{\Cotimes}{\mathscr{C}^{\otimes}}
\newcommand{\Ox}{\mathscr{O}^{\otimes}}
\newcommand{\Fin}{\mathscr{F} \text{in}_*}
\newcommand{\deltast}{\Delta_{\ast}}


%cocolon
\newcommand*\cocolon{%
        \nobreak
        \mskip6mu plus1mu
        \mathpunct{}%
        \nonscript
        \mkern-\thinmuskip
        {:}%
        \mskip2mu
        \relax
}

\tikzcdset{scale cd/.style={every label/.append style={scale=#1},
    cells={nodes={scale=#1}}}}

\numberwithin{equation}{section}
%bloedes Epsilon, Phi weg
\renewcommand{\epsilon}{\varepsilon}
% \renewcommand{\phi}{\varphi}

%Klickbare Links
\usepackage{hyperref}
